\documentclass{article}
\usepackage{graphicx}
\usepackage{amsfonts}
\usepackage{amssymb}
\usepackage{amsmath}
\usepackage{amsthm}
\usepackage{empheq}
\usepackage{graphicx}
\usepackage{color}
\usepackage{mathtools}
\usepackage[margin=1.5in]{geometry}
% \usepackage[polish]{babel}
\usepackage[utf8]{inputenc}
\usepackage[T1]{fontenc}

\title{Data Science}
\author{Filip Rabiega}

\theoremstyle{definition}
\newtheorem{theorem}{Theorem}[section]
\newtheorem{definition}{Definition}[section]
\newtheorem{proposition}{Proposition}[section]
\newtheorem{lemma}{Lemma}[section]
\newtheorem{fact}{Fact}[section]
\newtheorem{corollary}{Corollary}[section]
\newtheorem{remark}{Remark}[section]

\begin{document}

\maketitle
\section{Elementary facts}

\begin{definition}[Normal distribution]
  If \( X \sim N(\mu, \Sigma) \), then
  \[
    f_X (x_1, \dots, x_d) =
    \frac{1}{\sqrt{\left( 2 \pi \right) ^{d}  \text{det}(\Sigma)}}
    \exp\left(-\frac{1}{2}\left( x - \mu \right)^{T} \Sigma^{-1}
    \left( x - \mu \right)\right).
  \]
\end{definition}

\begin{definition}[\( \chi^{2} \) distribution]
  If \( X_{1}, X_{2}, \dots X_{n} \) are independent, standard normal
  variables, then the sum of their squares, \( Z = \sum_{i=1}^{n}
  X_{i}^{2} \), is distributed according to the \( \chi^{2} \)
  distribution with \( n \) degrees of freedom. We denote this by \(
  Z \sim \chi^{2} (n) \). Its probability density function is
  \[
    \varphi(t) =
    \left\{
      \begin{align}
        \frac{t^{\frac{n}{2}-1}e^{-\frac{n}{2}}}{2^{\frac{n}{2}}\Gamma(\frac{n}{2})}
        & \quad \text{if } t > 0, \\
        0 & \quad \text{else}.
      \end{align}
      \right.
    \]
  \end{definition}

  \begin{theorem}[Law of large numbers, version 1]
    If \( \operatorname{Var}(X _{i} ) < \infty \), then
    \[
      \mathbb{P} \left\{ | \bar{X} - \mathbb{E}X | \geq \varepsilon \right\}
      \leq \frac{\operatorname{Var}(X)}{n\varepsilon^{2}}
    \]
  \end{theorem}

  \begin{theorem}[Law of large numbers, version 2]
    If \( \operatorname{Var}(X) < \infty  \), then
    \[
      \sqrt{n} \left( \bar{X} - \mathbb{E}X \right)
      \xrightarrow{\mathcal{D} } N(0,\sigma ^{2} ).
    \]
  \end{theorem}

  \subsection{Key inequalities}

  \begin{theorem}[Markov's inequality]
    If \( X \) is a nonnegative random variable and \( a > 0 \), then
    \[
      \mathbb{P}\left\{ X \geq a \right\} \leq \frac{\mathbb{E}X}{a} .
    \]
  \end{theorem}

  \begin{remark}
    If \( X \) is an arbitrary random variable and \( a \) is an
    arbitrary real number then \( e^{X}, a>0 \), and so
    \[
      \mathbb{P}\left\{ \exp(X) \geq \exp(a)  \right\} \leq
      \frac{\mathbb{E}\exp(X) }{\exp(a) }
    \]
    by Markov's inequality.
  \end{remark}

  \begin{theorem}[Chebyshev's inequality]
    If \( a>0 \), then
    \[
      \mathbb{P}\left\{ | X - \mathbb{E}X | \geq a \right\} \leq
      \frac{\operatorname{Var}(X)}{a^2} .
    \]
  \end{theorem}

  \begin{theorem}[Jensen's inequality]
    Let \( \varphi: \mathbb{R} \rightarrow \mathbb{R} \) be convex. Then
    \[
      \varphi(\mathbb{E}X) \leq \mathbb{E}\varphi(X)
    \]
    for any random variable \( X \).
  \end{theorem}

  \begin{corollary}
    The function \( \varphi(x) = x^{p/q} \) is convex for \( p \geq q
    \geq 0 \), so
    \[
      \left\| X \right\|_{L^q} \leq \left\| X \right\|_{L^p}.
    \]
  \end{corollary}

  \begin{theorem}[Hoelder's inequality]
    For \( p,q \geq 1 \) such that \( \frac{p+q}{pq} = 1 \), the
    following holds:
    \[
      \left| \mathbb{E}XY \right| \leq \left\| X \right\|_{L^p} \left\|
      Y \right\|_{L^q}.
    \]
  \end{theorem}

  \begin{corollary}[Cauchy-Schwartz inequality]
    If \( X, Y \) are any random variables, then
    \[
      \left| \mathbb{E}XY \right| \leq \sqrt{\mathbb{E}X^{2} \mathbb{E}Y^{2}}.
    \]
  \end{corollary}

  \begin{proposition}
    If \( X \) is a nonnegative random variable, then
    \[
      \mathbb{E}X = \int_{0}^{\infty} \mathbb{P}\left\{ X > t \right\}dt.
    \]
  \end{proposition}

  \begin{proof}
    Fubini's theorem.
  \end{proof}

  \begin{proposition}[Gaussian tail bounds]
    If \( X \sim N(\mu, \sigma^{2}) \), then
    \[
      \mathbb{P}\left\{ X \geq \mu + t \right\} \leq e^{-t^{2} / \sigma^{2}}.
    \]
  \end{proposition}

  \begin{corollary}
    If \( X \sim N(0, \sigma^{2}) \), then
    \[
      \mathbb{P}\left\{ \left| X \right| \geq t \right\} \leq 2
      e^{-t^{2} / \sigma^{2}}.
    \]
  \end{corollary}

  \begin{definition}[Sub-Gaussian variables]
    A random variable \( X \) is sub-Gaussian if, for \( \mu =
    \mathbb{E}X \) and some \( K^{2} > 0 \),
    \[
      \mathbb{E}e^{\lambda(X - \mu)} \leq e^{K^{2} \lambda^{2} / 2}
    \]
    for all \( \lambda \in \mathbb{R} \).
  \end{definition}

  \begin{proposition}
    Assume that \( X \) is sub-Gaussian with parameter \( K^{2}
    \). Then for all \( t > 0 \),
    \[
      \mathbb{P}\left\{ \left| X - \mu \right| \geq t \right\} \leq
      2e^{-t^{2} / 2 K^{2}}.
    \]
  \end{proposition}

  \begin{theorem}[Hoeffding's inequality]
    Let \( X_i \) be pairwise independent Rademacher random variables
    and set \( a \in \mathbb{R}^d \).
    Then, for all \( t \geq 0 \),
    \[
      \mathbb{P}\left\{ \sum_{i=1}^{n} a_i X_i \geq t \right\} \leq
      \exp\left(-\frac{t^2}{2 \|a\|^{2}} \right).
    \]
  \end{theorem}

  \begin{proof}
    Set \( \lambda \geq 0 \). From Markov's inequality,
    \[ \mathbb{P}\left\{ \sum_{i=1}^{n} a_i X_i \geq t \right\} =
      \mathbb{P}\left\{ \exp\left( \lambda \sum_{i=1}^{n} a_i X_i \right) \geq
      \exp\left( \lambda t\right) \right\} \leq \exp\left(-\lambda t\right)
    \mathbb{E}\exp\left(\lambda \sum_{i=1}^{n} a_i X_i\right).\]
    Now we shall choose a value of \( \lambda \) that minimizes
    the right hand side. Notice that
    \[
      \mathbb{E} \exp\left(\lambda \sum_{i=1}^{n} a_i X_i\right) =
      \prod_{i=1}^{n} \mathbb{E}\exp\left(\lambda a_i X_i\right).
    \]
    Since \( X_i \) are all Radmacher-distributed, \( \mathbb{E}
    \exp\left( \lambda a_i X_i \right) = \cosh(\lambda a_i) \) by
    simple calculations.
    Now \( \cosh(x) \leq \exp\left(x^2 / 2\right) \), so \(
      \mathbb{E}\exp\left(\lambda a_i X_i\right) \leq \exp\left(
    \lambda^2 a_{i}^{2} / 2\right) \).
    Altogether
    \[
      \mathbb{P}\left\{ \sum_{i=1}^{n} a_i X_i \right\} \leq \exp\left(
      -\lambda t + \frac{\lambda^2}{2} \sum_{i=1}^{n} a_i ^2\right) =
      \exp\left(-\lambda t + \frac{\lambda^2}{2} \left\|a\right\|^2\right).
    \]
    The minimum is achieved for \( \lambda =
    \frac{t}{\left\|a\right\|^2} \). Then we get the desired inequality.
  \end{proof}

  \begin{remark}
    Hoeffding's inequality can be extended to cover many more types of
    random variables other than just Radmacher-distributed ones, for
    example bounded random variables.
  \end{remark}

  \begin{theorem}[Bernstein's inequality]
    Let \( X_1, \dots X_n \) be independent centered random variables
    satisfying \( \left| X_i \right| \leq a \) and \( \mathbb{E}
    X_{i}^{2} = \sigma^{2} \). Then
    \[
      \mathbb{P}\left\{ \left| \sum_{i=1}^{n} X_i \right| \geq t
      \right\} \leq 2 \exp\left( - \frac{t^{2}}{2n \sigma^2 +
      \frac{2}{3}at}\right).
    \]
  \end{theorem}

  \begin{theorem}[Chernoff's inequality]
    Let \( X_i \) be independent Bernoulli trials with \(
    \mathbb{P}\left\{ X_i = 1 \right\} = p_i \).
    Let \( S_n = \sum_{i=i}^{n} X_i \) and \( \mu = \mathbb{E}S_n \).
    Then
    \[
      \mathbb{P}\left\{ S_n \geq t \right\} \leq \exp(- \mu) \left(
      \frac{e \mu}{t}  \right) ^{t}
    \]
    for all \( t \geq \mu  \).
  \end{theorem}

  \begin{proof}
    Use the exact same strategy as in the proof of Hoeffding's inequality. Use
    \( 1 + x \leq \exp\left(x\right)\).
  \end{proof}

  \begin{corollary}
    For \( 0 \leq \delta \leq 1 \),
    \[
      \mathbb{P}\left\{ \left| S_n - \mu \right| \geq \delta \mu
      \right\} \leq 2 \exp\left(- \frac{\delta^2 \mu}{3}\right).
    \]
  \end{corollary}

  \begin{proposition}
    If \( Z = \sum_{i=1}^{n} X_{i}^{2} \sim \chi^{2}(n) \), then
    \[
      \mathbb{P}\left\{ \left| \frac{1}{n} \sum_{i=1}^{n} X_{i}^{2} -
      1 \right| \ge t \right\}
      \le
      \begin{cases}
        2e^{-nt^{2}/8} & \text{if } t \in (0,1), \\
        2e^{-nt/8} & \text{if } t \ge 1,
      \end{cases}
    \]
  \end{proposition}

  \begin{theorem}[Master Tail Bounds Theorem]
    Let \( X = \sum_{i=1}^{n} X_i \), where \( X_i \) are pairwise
    independent and have zero mean and variance at most \( \sigma ^{2}
    \). Let \( 0 \leq a \leq \sqrt{2}n \sigma ^{2} \). Assume that \( |
    \mathbb{E}X ^{s} _{i} | \leq \sigma ^{2} s! \) for \( s = 3, 4,
    \dots , \lfloor \frac{a^2}{4n \sigma^2} \rfloor \).
    Then
    \[
      \mathbb{P}\left\{ | X | \geq a \right\} \leq 3 \exp\left(-a^2 /
      12n \sigma^2 \right).
    \]
  \end{theorem}

  \subsection{Types of convergence}

  \begin{definition}[Convergence almost surely]
    We say that a sequence of variables \( X_n \) converges almost
    surely to \( X \) if
    \[
      \mathbb{P}\left\{ \omega : \lim_{n} X_{n}(\omega) =
      X(\omega) \right\} = 0.
    \]
  \end{definition}

  \begin{theorem}[Egorov]
    Let \( X_n \) be a sequence of random variables that converges
    almost surely to \( X \). For each \( \delta > 0 \) there is a set
    \( A_{\delta} \) such that \( \mathbb{P}\left\{ A_{\delta} \right\}
    < \delta \) and \( X_n \) converges to \( X \) uniformly on \(
    A_{\delta} \).
  \end{theorem}

  \begin{definition}[Convergence in probability]
    We say that a sequence of variables \( X_n \) converges in
    probability to \( X \) if for each \( \varepsilon > 0 \),
    \[
      \lim_{n \rightarrow \infty} \mathbb{P}\left\{ \left| X_n - X
      \right| > \varepsilon \right\} = 0.
    \]
  \end{definition}

  \begin{theorem}[First Borel-Cantelli lemma]
    If \( A_{n} \) are a sequence of events such that \(
    \sum_{n=1}^{\infty} \mathbb{P}\left\{ A_{n} \right\} < \infty \), then
    \[
      \mathbb{P}\left\{ \limsup_{n \rightarrow \infty} A_{n} \right\} = 0.
    \]
  \end{theorem}

  \begin{theorem}[Second Borel-Cantelli lemma]
    If \( A_{n} \) are pairwise independent events and \( \sum_{n=1}^{\infty}
    \mathbb{P}\left\{ A_{n} \right\} = \infty \), then
    \[
      \mathbb{P}\left\{ \limsup_{n \rightarrow \infty} A_{n} \right\} = 1.
    \]
  \end{theorem}

  \begin{lemma}
    If \( X_n \rightarrow X \) in probability, then there is a
    subsequence \( X_{n_{k}} \) that converges almost surely to \( X \).
  \end{lemma}

  \begin{lemma}
    \( X_{n} \rightarrow X \) almost surely iff \( \sup_{k
    \geq n} \left| X_{k} - X \right| \rightarrow 0 \) in probability.
  \end{lemma}

  \begin{corollary}
    Almost sure convergence implies convergence in probability.
  \end{corollary}

  \begin{theorem}[Dominated Convergence Theorem]
    Let \( X_n \) be a sequence that converges in probability to \( X
    \). Suppose that \( \mathbb{E} \sup_{n} \left| X_{n} \right| <
    \infty \). Then \( \mathbb{E} \left| X_{n} - X \right| \rightarrow 0\).
  \end{theorem}

  \begin{definition}[Convergence in \( L_{p} \)]
    We say that \( X_{n} \) converges to \( X \) in \( L_{p} \) for \(
    1 \leq p < \infty \) if
    \[
      \lim_{n} \mathbb{E} \left| X_{n} - X \right|^{p} = 0.
    \]
  \end{definition}

  \begin{fact}
    Convergence in \( L_{p} \) implies convergence in probability. On
    the other hand, it is independent of convergence almost surely,
    i.e. one does not imply the other in any direction.
  \end{fact}

  \begin{definition}[Convergence in distribution]
    We say that a sequence \( X_{n} \) converges in distribution to \(
    X \) if for each continuous bounded function \( f: \mathbb{R}
    \rightarrow \mathbb{R} \), the convergence
    \[
      \lim_{n} \mathbb{E}f(X_{n}) = \mathbb{E}f(X)
    \]
    takes place.
  \end{definition}

  \begin{proposition}
    A sequence \( X_{n} \) converges in distribution to \( X \) iff for
    each continuity point of the CDF of \( X \),
    \[
      \lim_{n} \mathbb{P}\left\{ X_{n} \leq x \right\} =
      \mathbb{P}\left\{ X \leq x \right\}.
    \]
  \end{proposition}

  \begin{proposition}
    If a sequence \( X_{n} \) converges in probability to \( X \), then
    it converges in distribution to \( X \).
  \end{proposition}

  \begin{theorem}[Continuous Mapping Theorem]
    If \( X_{n} \) converges in distribution to \( X \) and \( g:
    \mathbb{R} \rightarrow \mathbb{R} \) is continuous, then \(
    g(X_{n}) \) converges in distribution to \( g(X) \).
  \end{theorem}

  \section{Geometry in high-dimensional spaces}
  \begin{theorem}
    The volume of a unit ball in a \( d \)-dimensional space is
    \[
      \text{Vol}(\mathbb{B}^d) = \frac{\pi ^{d/2} }{\frac{d}{2} \Gamma
      (\frac{d}{2} )}.
    \]
  \end{theorem}

  \begin{proof}
    Induction with respect to \( d \).
  \end{proof}

  \begin{corollary}
    For big \( d \), this volume is close to \( 0 \).
    Also, the volume of the ball is concentrated near the boundary.
  \end{corollary}

  \begin{fact}
    If \( X \sim N(0, I_d) \), then
    \[
      \begin{aligned}
        \mathbb{E} \|X\|^{2} &= d, \\
        \operatorname{Var}(\|X\|^2) &= 2d.
      \end{aligned}
    \]
  \end{fact}

  \begin{fact}
    If \( X \sim N(0, I_d) \), then
    \[
      \begin{aligned}
        \left| \mathbb{E}\left( \|X\| - \sqrt{d} \right) \right|
        &\leq 1/\sqrt{d}, \\
        \operatorname{Var}(\|X\|) &\leq 2.
      \end{aligned}
    \]
  \end{fact}

  \begin{remark}
    The facts above tell us that the Gassian variables in \(
    \mathbb{R}^d \) are concentrated around a sphere with radius \(
    \sqrt{d} \).
  \end{remark}

  \begin{fact}
    If \( X, Y \sim N(0, I_d) \) pairwise independent and \( \xi \in
    \mathbb{R}^d \), then
    \[
      \begin{align}
        \mathbb{E}\langle X, Y \rangle &= 0, \\
        \operatorname{Var}(\langle X, Y \rangle) &= d, \\
        \mathbb{E}\langle X, \xi \rangle &= 0, \\
        \operatorname{Var}(\langle X, \xi \rangle) &= \|\xi\|^{2}.
      \end{align}
    \]
  \end{fact}

  \begin{theorem}[Ortogonality of Gaussian variables]
    If \( X \sim N(0, I_d) \) and \( 0 \leq \varepsilon \), then
    \[
      \mathbb{P}\left\{ \left| \left\langle \frac{X}{\|X\|}, \frac{Y}{\|Y\|}
      \right\rangle \right| \leq \varepsilon \right\} \geq 1 - \frac{2 /
      \varepsilon + 7}{\sqrt{d}}.
    \]
  \end{theorem}

  \begin{theorem}[Gaussian Annulus Theorem]
    Let \( X \sim N(0, I_d) \) and \( \varepsilon \leq \sqrt{d} \). Then
    \[
      \mathbb{P}\left\{ \left| \left\|X\right\| - \sqrt{d} \right| \leq
      \varepsilon \right\} \geq 1 - 3 \exp\left( -c \varepsilon^2\right),
    \]
    where \( c > 0 \) is some constant.
  \end{theorem}

  \begin{proof}
    Let \( r = \left\| X \right\| \). If \( \left| r - \sqrt{d} \right|
    \geq \varepsilon \), then \( \left| r^2 - d \right| \geq
      \varepsilon \left( r + \sqrt{d} \right) \geq \varepsilon \sqrt{d}
    \). So it suffices to bound the probability that \( \left| r -
    \sqrt{d} \right| \geq \varepsilon \sqrt{d} \).

    Rewrite \( r^2 - d = (X_{1}^{2} - 1) + \dots + (X_{d}^{2} - 1)\)
    and set \( Y_i = X_{i}^{2} - 1 \). We want to bound the probability
    that \( \left| \sum_{i=1}^{d} Y_i \right| \geq \varepsilon \sqrt{d}
    \). Notice that \( \mathbb{E} Y_i = 0\). To apply the master
    theorem we have to bound the moments of \( Y_i \). But that is
    simple calculus; in the end we get \( \left| \mathbb{E} Y^{s}
    \right| \leq 2^{s}s! \).
    Variance of \( Y_i \) is still to high, so we set \( Z_i = Y_i / 2
    \), then \( \operatorname{Var}(Z_i) \leq 2 \) and \( \left|
    \mathbb{E}Z_{i}^{s} \right| \leq 2s! \). Now we can apply the
    Master theorem for \( \sigma^2 = 2 \) and \( n = d \), which yields
    the result.
  \end{proof}

  \begin{theorem}[Random Projection Theorem]
    Set \( v \in \mathbb{R}^d \) and let \( f(v) = Uv \), where \( U
    \in M_{n \times k}(\mathbb{R}) \) the coefficients of \( U \) are
    standard and pairwise independent Gaussians. Then there exists a
    constant \( c > 0 \) such that, for \( \varepsilon \in (0,1) \),
    \[
      \mathbb{P}\left\{ \left| \left\| f(v) \right\| - \sqrt{k} \left\| v
        \right\| \right| \geq \varepsilon \sqrt{k} \left\| v \right\|
      \right\}  \leq 3 \exp\left(-ck \varepsilon^2\right).
    \]
  \end{theorem}

  \begin{proof}
    By scaling the inner inequality by \( \left\| v \right\| \), we can
    assume that \( \left\| v \right\| = 1 \). Since the sum of normal
    variables is normal, \( \left\langle u_i, v \right\rangle \) is
    Gaussian with zero mean and variance equal to
    \[
      \operatorname{Var}(\left\langle u_i, v \right\rangle) =
      \sum_{j=1}^{d} v_{j}^{2} u_{ij} = \sum_{j=1}^{d} v_{j}^{2} = 1.
    \]
    Since \( \left\langle u_{i}, v \right\rangle \) for \( i=1, \dots k \) are
    independent, Gaussian and have unit variance, the desired property
    follows from the Gaussian Annulus Theorem.
  \end{proof}

  \begin{theorem}[Johnson-Lindenstrauss Lemma]
    For any \( 0 < \varepsilon < 1 \) and any positive integer \( n \),
    let \( k \geq \frac{3}{c \varepsilon^2} \log n \) with \( c \) as
    in the Gaussian Annulus Theorem. For any set of \( n \) points in
    \( \mathbb{R}^d \), the random projection \( f(v) = Uv \), where
    \( U \) is as before, has the property that for all \( v_{i}, v_{j} \),
    \[
      \mathbb{P}\left\{ (1-\varepsilon) \sqrt{k} \left\|v_{i} -
        v_{j}\right\| \leq \left\| f(v_{i}) - f(v_{j})\right\| \leq
      (1+\varepsilon) \sqrt{k} \left\|v_{i} - v_{j}\right\| \right\}
      \geq 1- 3/2n.
    \]
  \end{theorem}

  \begin{proof}
    Simply use the previous result with an union bound with \(
    \binom{n}{2} < n^{2}/2 \).
  \end{proof}

  \begin{fact}
    Let \( X \) be a random vector and let \( \Sigma = \mathbb{E}XX^T
    \) be its covariance matrix. Then
    \begin{itemize}
      \item for all \( v \in \mathbb{R}^d \), \( \mathbb{E}\left\langle
        X, v \right\rangle^2 = v^T \Sigma v \).
      \item \( \mathbb{E}\left\|X\right\|^2 = \operatorname{tr}(\Sigma) \),
      \item If \( Y \) is an independent copy of \( X \), then \(
          \mathbb{E}\left\langle X, Y \right\rangle^2 =
        \left\|\Sigma\right\|^{2}_{F} \), where \( \left\|\cdot\right\|_F \)
        stands for the Frobenius norm.
    \end{itemize}
  \end{fact}

  \section{SVD Decomposition}

  SVD decomposition is a way to project a set of points onto a small
  subspace in such a way that the sum of squares of distances from the
  original points to the projections is minimal.

  \begin{definition}[Singular Value Decomposition]
    SVD is a decomposition of a matrix \( A \in M_{n \times
    d}(\mathbb{R}) \), where we think of rows of \( A \) as points in
    a \( d \)-dimensional space, is the factorization
    \[
      A = U \Sigma V^{T},
    \]
    where the columns of \( U, V \) are orthonormal and the matrix \(
    \Sigma \) is diagonal with positive entries.
  \end{definition}

  \begin{remark}
    In the above definition the columns of \( V \) are the vectors
    defining the best-fitting subspace.
  \end{remark}

  \begin{remark}
    If \( \operatorname{rank}(A) = r \), the decomposition above can
    be rewritten as
    \[
      A = \sum_{i}^{r} \sigma_{i}u_{i}v_{i}^{T},
    \]
    where \( \sigma_{i} > 0 \) for each \( 1 \leq i \leq r\).
    Here each \( u_{i}v_{i}^{T} \) is a \( n \times d \) matrix of rank \( 1 \).
  \end{remark}

  \begin{remark}
    If \( A \) is symmetric, then the singular values are the absolute
    values of eigenvalues.
  \end{remark}

  \begin{lemma}
    If \( A = U \Sigma V^{T} \) is the SVD, then
    \[
      Av_{i} = \sigma_{i} u_{i}, \; A^{T}u_{i} = \sigma_{i} v_{i}.
    \]
  \end{lemma}

  \begin{remark}
    The greedy algorithm works when calculating \( V \). Simply notice that
    \( v_{1} = \arg\max_{\left\|v\right\| = 1} \left\|Av\right\| \) and
    generate subsequent vectors from subspaces orthogonal to the ones
    spanned by the already found columns of \( V \).
    The fact that this approach works can be proven by simple induction.
  \end{remark}

  \begin{remark}
    If \( A = U \Sigma V^{T} \), then the \( i \)-th singular value is \(
    \sigma_i (A) = \left\|Av_{i}\right\| \).
  \end{remark}

  \begin{lemma}[Frobenius norm vs. singular values]
    \[
      \left\|A\right\|_{F} = \sqrt{\sum_{i}^{} \sigma_i (A)^{2}}.
    \]
  \end{lemma}

  \begin{proof}
    \[
      \left\| A \right\|^{2}_{F} = \sum_{i,j} a_{ij}^{2} =
      \operatorname{tr}(A^{T}A) = \operatorname{tr}(V \Sigma^{T} U^{T}
      U \Sigma V^{T}) = \operatorname{tr}(\Sigma^{T} \Sigma) =
      \sum_{i=1}^{r} \sigma_{i}^{2}.
    \]
  \end{proof}

  \begin{definition}
    \[
      u_i = \frac{1}{\sigma_{i}(A)} Av_{i}.
    \]
  \end{definition}

  \begin{remark}
    Notice that \( \sigma_{i} u_{i}v_{i}^{T} \) is a \( 1
    \)-dimensional matrix whose rows are the \( v_{i} \)-components
    of \( A \).
  \end{remark}

  \begin{remark}
    \[
      \left\| Ax \right\|^{2} = \sum_{i=1}^{r} \sigma_{i}^{2}
      \left\langle v_{i}, x \right\rangle^{2} \leq \sigma_{1}^{2}
      \sum_{i=1}^{r} \left\langle v_{i}, x \right\rangle^{2} \leq
      \sigma_{1}^{2} \sum_{i=1}^{n} \left\langle v_{i}, x
      \right\rangle^{2} \leq \sigma_{1}^{2} \left\| x \right\|^{2}.
    \]
  \end{remark}

  \begin{corollary}
    \[
      \left\| A \right\|_{2} = \sigma_{1},
    \]
    where \( \left\| A \right\|_{2} = \sup_{\left\| x \right\| = 1}
    \left\| Ax \right\| \) by selecting \( x = v_{1}\).
  \end{corollary}

  \begin{lemma}[Best rank-\( k \) approximations]
    If \( A = \sum_{i=1}^{r} \sigma_{i}u_{i}v_{i}^{T} \) is the SVD, then
    \[
      A_k = \sum_{i=1}^{k} \sigma_{i} u_{i}v_{i}^{T}
    \]
    is the best rank-\( k \) approximation of \( A \) in the Frobenius norm.
  \end{lemma}

  \begin{lemma}
    \[
      \begin{aligned}
        \left\|A - A_{k}\right\|_{2}^{2} &= \sigma_{k+1}^{2}, \\
        \left\|A - A_{k}\right\|_{F}^{2} &= \sum_{i=k+1}^{r} \sigma_{i}^{2}. \\
      \end{aligned}
    \]
  \end{lemma}

  \begin{fact}
    The centroid of a set of data points must pass through the
    best-fitting line.
  \end{fact}

  \begin{lemma}[Principal Component Analysis]
    Let \( X \) be a random vector with \( \lambda_1 \geq \lambda_2
    \geq \dots \geq \lambda_n \geq 0  \) as the eigenvalues of its
    covariance matrix. Then
    \[
      \lambda_k = \max_{v \perp \left\{ v_1, \dots, v_{k-1} \right\},
      \left\|v\right\| = 1} \operatorname{Var}(\left\langle X,
      v \right\rangle).
    \]
  \end{lemma}

  \section{Machine Learning}

  \subsection{Perceptron}

  \begin{definition}[Linear separability]
    We say that a set of vectors \( W \subseteq \mathbb{R}^d \times
    \left\{ -1, 1 \right\} \) is \emph{linearly separable with
    threshhold \( \gamma \)}, if there is a vector \( v \in \mathbb{R}
    \) such that for all \( \left( x,y \right) \in W \),
    \[
      \frac{y \left\langle v, x \right\rangle}{\left\|v\right\|} \geq \gamma.
    \]
  \end{definition}

  \begin{remark}
    As is widely known, \( \frac{\left| \left\langle v, x \right\rangle
    \right|}{\left\|v\right\|} \) is the distance from \( x \) to the
    hyperplane orthogonal to \( v \).
  \end{remark}

  \begin{definition}[Affine linear classificator]
    For \( W = \left\{ (x_i, y_i) \; : \; i = 1, 2, \dots, n \right\}
    \subseteq \mathbb{R}^d \times \left\{ -1, 1 \right\} \), an
    \emph{affine linear classificator} is a function
    \[
      h: \mathbb{R}^d \rightarrow \mathbb{R},\; h(x) =
      \operatorname{sgn}(\left\langle w, x \right\rangle + b),
    \]
    where \( h(x_i) = y_i \) for all \( i \).
  \end{definition}

  \begin{remark}
    We will assume that \( w = \left( w_1, w_2, \dots, w_n, b \right)
    \) and \( x = \left( x_1, x_2, \dots, x_n, 1 \right) \).
  \end{remark}

  \end{document}
